\chapter{Đặt vấn đề}

\section{Giới thiệu về bài toán}

Trong bối cảnh chuyển đổi số mạnh mẽ của ngành giáo dục, nhu cầu số hóa tài liệu
giảng dạy đang gia tăng đột biến. Một trong những nút thắt lớn nhất hiện nay là quy
trình chuyển đổi các đề thi, bài tập từ dạng tài liệu giấy hoặc ảnh chụp sang
định dạng dữ liệu số hóa để lưu trữ và tái sử dụng trên các hệ thống quản lý học
tập.

Hiện tại, quy trình này vẫn phụ thuộc nhiều vào thao tác thủ công hoặc các công cụ
OCR truyền thống (như Tesseract). Tuy nhiên, các công cụ này thường gặp khó khăn
khi xử lý các tài liệu có cấu trúc phức tạp hoặc chất lượng ảnh thấp. Sự xuất
hiện của các Mô hình Ngôn ngữ Lớn Đa phương thức như GPT-4o hay Gemini Pro đã mở
ra hướng giải quyết mới với độ chính xác vượt trội. Mặc dù vậy, rào cản về chi
phí vận hành, độ trễ và lo ngại về bảo mật dữ liệu đã ngăn cản việc triển khai rộng
rãi các mô hình khổng lồ này trong các ứng dụng giáo dục thực tế.

Xuất phát từ thực tiễn đó, việc nghiên cứu ứng dụng các Mô hình Ngôn ngữ Lớn cỡ
nhỏ (Small Large Language Models - SLLMs) với số lượng tham số dưới 10 tỷ (như Qwen-VL,
Llama-3-Vision, Nemotron) trở thành một hướng đi đầy tiềm năng. Các mô hình này hứa
hẹn mang lại sự cân bằng tối ưu giữa hiệu năng xử lý, chi phí vận hành và khả
năng triển khai linh hoạt. Đề tài "Nghiên Cứu, Đánh Giá Hiệu Năng sLLMs Trong Bài
Toán Sinh Câu Hỏi Trắc Nghiệm Từ Ảnh và Xây Dựng Ứng Dụng Thực Nghiệm" được thực
hiện nhằm kiểm chứng giả thuyết này.

\section{Các nghiên cứu liên quan}

\section{Hướng tiếp cận và đóng góp của khoá luận}

Mục tiêu tổng quát của đề tài là xây dựng một quy trình xử lý khép kín sử dụng
các mô hình SLLMs để tự động hóa việc trích xuất nội dung từ ảnh chụp đề thi và
sinh ra các câu hỏi trắc nghiệm tương đương, đảm bảo chất lượng học thuật và tối
ưu chi phí.

Các mục tiêu cụ thể bao gồm:

- \textbf{Đánh giá thực nghiệm}: So sánh hiệu năng của các mô hình SLLMs mã nguồn
mở với các mô hình thương mại trên các khía cạnh: khả năng nhận dạng văn bản, khả
năng suy luận logic và chất lượng sinh văn bản.

- \textbf{Tối ưu hóa dữ liệu}: Đề xuất quy trình sinh dữ liệu giả lập sử dụng kỹ
thuật Browser Engineering để khắc phục tình trạng thiếu hụt dữ liệu huấn luyện OCR
chất lượng cao.

- \textbf{Giải quyết bài toán kỹ thuật}: Xây dựng chiến lược Prompting và định dạng
đầu ra để khắc phục các hạn chế cố hữu của mô hình nhỏ (như ảo giác, lỗi cú pháp
JSON).

- \textbf{Xây dựng ứng dụng minh họa}: Phát triển ứng dụng Web "Openlet" tích
hợp các kết quả nghiên cứu, chứng minh tính khả thi khi triển khai thực tế.